\slajdid{08-s053}
\begin{frame}[label=08-s053]
\gusttekst\setlength{\parskip}{2pt}\setlength{\abovedisplayskip}{3pt}\setlength{\belowdisplayskip}{3pt}
\[p(t)=p(\infty)+\bigl(p(t_0^+)-p(\infty)\bigr)e^{-\frac{t-t_0}{\tau}}\qquad za\ t\geq t_0\]
Znači nećemo pisati i rešavati diferencijalne jednačine. Krenućemo od gotovog rešenja. Na nama je samo da odredimo
{\color{DEcrvena}\[\tau\qquad p(t_0^+)\qquad p(\infty)\]}
1. $\tau$\par
Određujemo kao proizvod C i otpornosti koju vidi taj kondenzator tako što sve nezavisne naponske izvore kratko spojimo, nezavisne strujne izvore uklonimo iz kola, ostavimo otvorene veze, pri čemu zavisni izvori ostaju.\par
2. $p(t_0^+)$\par
Određujemo na osnovu saznanja da je u pitanju realno kolo i da napon na kondenzatoru ne može trenutno da se promeni.
\[u_c(t_0^+)=u_c(t_0^-)\]
3. $p(\infty)$\par
Određujemo na osnovu saznanja da su u beskonačnosti završeni prelazni procesi odnosno da nema promene napona na kondenzatoru. A uslov da nema promene napona na kondenzatoru je $i_c(\infty)=0$ i zamenjujemo u gotovo rešenje
\[p(t)=p(\infty)+\bigl(p(t_0^+)-p(\infty)\bigr)e^{-\frac{t-t_0}{\tau}}\qquad za\ t\geq t_0\]
\end{frame}
